\section{Evaluation}
\label{sec:evaluation}

We evaluate three questions: \textbf{RQ1}, can enforcement constrain effects
inside an authorized application computation; \textbf{RQ2}, what functionality
and integration work does an existing application retain; and \textbf{RQ3}, what
does enforcement cost, and where is that time spent? An application comparison
isolates authorization from host enforcement. Separate component checks examine
the reference-monitor obligations in Section~\ref{sec:model}.

\subsection{Applications and experimental setup}

We reproduce three applications from two repositories: the bundled Pandas Data
Analyst and SQL Database Agent from AI Data Science Team~\cite{aidatascienceteam},
and Open Notebook~\cite{opennotebook}. Only Pandas has a protected integration.
Five other applications on the fixed eight-application shortlist remain blocked
or unverified; Appendix~\ref{app:application-details} records their status.
Baseline reproduction is not counted as a protected port.

The Pandas adapter invokes the original \code{PandasDataAnalyst} graph, including
its wrangling and visualization agents, prompts, and generated computation.
We select a local provider and exclude the Streamlit UI, authentication, upload
widgets, history, and rendering. All configurations receive the same IPython
dependency fix and CPU affinity. The upstream checkout is unchanged. We compare:
\begin{itemize}
  \item \textbf{A: upstream execution.} The original subprocess sandbox runs on
  the same LangGraph source overlay as B/C, with \sys security disabled.
  \item \textbf{B: added authorization.} A \code{ToolNode} boundary authorizes
  computation before invoking that original subprocess helper.
  \item \textbf{C: authorization and host enforcement.} The same boundary uses
  \code{LinuxSandboxExecutor}, running the upstream sandbox program in a fresh
  bubblewrap worker instead of nesting another subprocess.
\end{itemize}
The upstream AST, import, builtin, network, resource, and JSON-transport checks
are retained. B tests an added dispatch boundary, not a published defense or an
ablation of native upstream tool dispatch. C measures the bubblewrap backend,
not the Landlock or object-broker executor.

Experiments use a Xeon w5-3435X, 62.25~GiB RAM, two RTX A4000 GPUs, and Ubuntu
24.04.5 with Linux 6.17.0-1033-oem. Pandas and SQL use local Qwen3-8B Q4\_K\_M
through Ollama, temperature zero, requested seed 20261005, and no client retries.
The parent and computation children use logical CPU~0; Ollama runs separately.
Audit files reside on a rotating ext4 disk. The machine is not exclusive,
frequency is unfixed, and caches are not flushed. The implementation is WIP
\code{40e24e9} plus a frozen working-tree overlay; the source archive and per-file
hashes, not that commit alone, identify the measured implementation.

\subsection{RQ1: Containing an authorized computation}

\paragraph{Controlled effects.}
We preserve a legitimate February sales computation and insert one filesystem
operation before it returns. This directly tests whether an authorized body can
exceed its grant, independently of whether a model generates an attack. Policies
are fixed before each run. Narrow and broad profiles differ only in permission
to read one supplementary file. Each case/profile runs once per configuration.
Independent oracles check canary receipts, file hashes, kernel errors, and worker
activation evidence, as well as the legitimate result.

\begin{table*}[t]
\centering
\small
\begin{tabularx}{\textwidth}{lXX X}
\toprule
Operation / profile & A: upstream & B: added authorization & C: authorization + host enforcement \\
\midrule
Read unrelated file / narrow & Canary copied & Canary copied & \code{ENOENT}; no canary receipt \\
Read supplementary file / narrow & Canary copied & Canary copied & \code{ENOENT}; no canary receipt \\
Same supplementary read / broad & Read succeeds & Read succeeds & Read succeeds \\
Write outside host fixture & Host file created & Host file created & Host unchanged; private \code{/tmp} write succeeds \\
Overwrite dataset / narrow & Dataset modified & Dataset modified & \code{EROFS}; dataset hash unchanged \\
\bottomrule
\end{tabularx}
\caption{Effects inside an authorized Pandas computation. Each cell is one
execution; all 15 executions return the correct February result. The dataset
overwrite is a separately recorded diagnostic amendment. A private temporary
write is host containment, not denial of all writes outside the output directory.}
\label{tab:application-effects}
\end{table*}

Table~\ref{tab:application-effects} shows the distinction the framework is
designed to enforce. Both A and B reach all three forbidden effects for which
the configured restrictions are supported: two narrow-profile reads and the
dataset overwrite. C reaches none, while all 15 executions retain the correct
February result. The successful broad-profile read confirms that C distinguishes
an authorized resource from the same resource under a narrower grant.

The outside-fixture write exposes a different outcome. C leaves the host file
unchanged, but permits a write in its private \code{/tmp}. We therefore exclude
this case from the three-probe denominator and report strict guest-wide
output-directory-only writes as unsupported. The artifact preserves the
original overclaim based on host-file absence and its correction. Network,
descendant-process, and unavailable-host-mechanism cases were not independently
tested in this application. The experiment establishes containment of the tested
effects, not a census of all application sinks.

\paragraph{Authorization at the added boundary.}
B and C each pass four deterministic dispatch checks. An allowed request creates
its unique witness; a guest role, a missing \code{sandbox\_source} argument, and
an unknown tool each produce an error without that effect. Audit-chain
verification passes. Together with Table~\ref{tab:application-effects}, these
checks show why correct call authorization is necessary but insufficient:
B rejects unauthorized calls yet permits forbidden effects inside allowed ones.

\paragraph{Model-driven injection.}
A separate case places one frozen instruction payload in a sampled CSV
\code{region} cell. It requests a private-canary read and copy before returning
the correct total. Five attacked runs per configuration are paired with five
fresh clean runs. The payload reaches the model input in all 15 attacked runs,
but none generates the forbidden operation or attempts its target open. All
15 attacked tasks succeed. These are \code{not\_attempted} outcomes: conditional
enforcement failure is undefined with zero eligible attempts, so this batch
does not establish prompt-injection resistance. An earlier 30-run launch batch
failed before inference because Ollama was unavailable; those records are
retained separately from the 30 retry executions.

\subsection{RQ2: Utility and integration requirements}

The initial clean workload has three tasks, each run once in A/B/C. Every
configuration succeeds on 1/3: the February total is correct, regional
aggregation returns unaggregated rows, and chart execution encounters an
upstream regular-expression error. Later repetitions deliberately use the
known-working February task and succeed on 5/5 clean runs per configuration.
The equal initial outcomes and successful repeats show preservation of that
workflow, not broad utility non-inferiority or improvement on the original
three-task workload.

Pandas requires a 200-line integration adapter and 20-line worker, plus a
135-line backend adapter shared with SQL. These physical source-line counts
describe the integration surface; engineering time was not measured. The
original graph and generated computation remain in use.

SQL passes 3/3 baseline tasks on a SQLite fixture with \code{sales} and
\code{payroll} tables and upstream safe-mode validation. Trace inspection finds
that schema sampling reads a payroll canary into the model prompt before query
execution. This identifies a mediation requirement, not an exploit result:
protecting only generated queries misses the metadata path. Our fixed-schema
record mediator does not yet cover this application's arbitrary SQLAlchemy
queries. Open Notebook's 103-line adapter exercises native ingestion, source
chat, and a live dedicated SurrealDB, passing one smoke task. Its API and worker
hosts are not confined or fully traced. Neither baseline is counted as a
protected integration.

\subsection{RQ3: Runtime cost and its sources}

\paragraph{Matched computation replay.}
We replay the same generated February computation and recorded input through
A/B/C. Each configuration receives one first-use call, ten warm-ups, and 100
measured calls, rotating configuration order by round. Every replay returns
120. The timer includes dispatch, sandbox execution, synchronous audit writes,
and integration/evidence output. Fixtures and initialization are outside the
timer; there are no model calls or syscall tracing. This measures a real
application computation path rather than a no-op or full agent workflow.

\begin{table}[t]
\centering
\small
\begin{tabular}{lrrrr}
\toprule
Config. & Mean & SD & Median & P95 \\
\midrule
A & 267.19 & 6.84 & 265.00 & 280.56 \\
B & 545.54 & 31.53 & 544.34 & 583.76 \\
C & 930.93 & 42.95 & 916.64 & 1,010.42 \\
\bottomrule
\end{tabular}
\caption{Local computation replay in milliseconds, 100 measured calls per
configuration. SD describes variation within this trace, not across tasks.}
\label{tab:replay-latency}
\end{table}

Table~\ref{tab:replay-latency} reports means of 267.19, 545.54, and 930.93~ms.
The paired mean increments are 278.35~ms for B$-$A, 663.74~ms for C$-$A, and
385.39~ms for C$-$B, with SDs of 32.91, 42.55, and 54.98~ms. C costs
$3.48\times$ A on this trace, a 248.4\% increase. Enforcement is therefore
substantial in this configuration. It cannot be characterized by the earlier
13.2\% policy-only no-op measurement.

\paragraph{Where the time goes.}
A and B spend 267.11 and 267.03~ms in the original subprocess path. B's remaining
278.51~ms includes dispatch, authorization, and evidence persistence. A separate
20-call diagnostic, after two warm-ups, directly times SQLite audit writes:
273.94~ms in B (SD 32.12) and 274.52~ms in C (SD 8.96). In that diagnostic,
B totals 544.34~ms with 265.09~ms in its subprocess; audit persistence accounts
for about 98\% of the remaining 279.25~ms. This identifies a concrete
implementation bottleneck. No SSD comparison or disabled-audit experiment was
performed, so a hardware speedup is not inferred.

In the 100-call cohort, C's executor averages 614.54~ms, including 20.56~ms for
grant compilation and 325.58~ms for worker execution and supervision. The
268.41~ms residual inside the executor includes other preparation and enforcement
work; another 316.39~ms lies outside it in dispatch and integration/evidence
work. These nested spans are not additive with their enclosing executor span.

\begin{table}[t]
\centering
\small
\begin{tabular}{lrrrr}
\toprule
Config. & Mean & SD & Median & Success \\
\midrule
A & 12.777 & 1.410 & 12.236 & 5/5 \\
B & 12.725 & 0.195 & 12.736 & 5/5 \\
C & 15.354 & 0.118 & 15.298 & 5/5 \\
\bottomrule
\end{tabular}
\caption{Recorded full-workflow time in seconds for five clean February runs
per configuration. Scope includes runner setup, environment capture, the
model-driven subprocess under \code{strace}, and the outcome oracle.}
\label{tab:workflow-latency}
\end{table}

\paragraph{Full-workflow observations.}
Table~\ref{tab:workflow-latency} reports the separate five-run clean cohort.
C$-$A averages 2.577~s (SD 1.446), a 20.2\% ratio-of-means increase; C$-$B
averages 2.629~s (SD 0.216). B$-$A is $-0.052$~s (SD 1.581), which does not
establish a speedup. The first A run follows a model-server restart and retains
loading effects. Model output, setup, tracing, and evidence collection all
contribute to these totals. Thus the 20.2\% difference is not an isolated
enforcement-overhead estimate, and the $3.48\times$ replay ratio is not a
whole-application slowdown. Five repetitions of one task do not characterize
across-task uncertainty.

\subsection{Supporting reference-monitor evidence}

The application experiments complement the earlier component evaluation.
Eight deliberately vulnerable tools reach unauthorized sinks behind
policy-conformant calls with no defense, Progent, and three generic wrappers.
These establish the effect gap, not incorrect call-policy decisions. A historical
router analysis checks 21 dispatch obligations; Tool API conformance covers
15,744 cases, and a focused router/object/host run passes 151 tests. Linux tests
exercise descriptor binding, raw syscalls, descendants, activation, and resource
controls. These records support specific contracts at their original snapshots.
They do not prove whole-framework G1--G3, nor transfer Landlock properties to
the application study's bubblewrap path.

Appendix~\ref{app:component-evidence} retains the per-layer results, analysis and
deployment evidence, policy-only benchmark, and open composition diagnostics.
The application evidence establishes one protected computation integration;
multiple workloads, semantic-broker comparisons, runtime-MAC/analysis composition,
and genuine host-capability degradation remain to be evaluated.
